\section{Ideals}

\begin{definition}
A subring  $I$ of a ring $R$ is a two-sided ideal) if for all $x\in R$ and $i\in I, xi \mbox{ and } ix$ are both in $I$.
\end{definition}

Note that because $R$ may not be a commutative ring we must have both $xi$ and $ix$ in $I$. 

\begin{theorem}
If $R$ has a unity then 
\[ \{ xi: x\in R, i \in I \} \subseteq I \implies \{xi\} = I \mbox{ and similarly for } \{ ix \}.\]
\end{theorem}
\begin{proof}
$ \{ xi \} \supseteq \{ 1i \} = I $
\end{proof}

{\bf Quotient Rings }
\begin{theorem}
If $I$ is an ideal of a ring $R$ the cosets $x + I$ form a ring under the sum $ (x + I) + (y + I) = (x+y) + I$ and products
$(x + I ) (y + I) = xy + I$.
\end{theorem}

The ring of cosets is called the {\twelveit Quotient Ring of I in R} and it is written $R/I$.  The mapping $a \goes a' = a + I$ for $a \in R$, which carries each element of $R$ into the coset 
containing it is a homomorphism of $R$ into the quotient ring $R/I$ and the elements mapped on zero in this homomorphism are exactly the elements of the ideal $I$. In the homomorphic ring, the zero is the coset $0 + I$, hence the elements of $I$ are mapped onto the zero coset.  

\begin{theorem}
If a homomorphism $H$ maps $R$ on $R'$ and exactly the elements of $I$ on zero, then $R'$ is isomorphic to the quotient ring $R/I$. 
\end{theorem}

If $b$ is an element of a commutative ring with unity, the set $[b]$ of all multiples $xb$ of $b$, for any $x\in R$, is an ideal known as a {\twelveit principal ideal}; it is the smallest ideal of $R$ containing $b$. \\

Inclusion between ideals is closely related to divisibility between numbers. In the ring $J$ it integers $n|m$ means that $m = an$, for some integer $a$ ($n$ divides $m$). Hence, every multiple of $m$ is a multiple of $n$. The multiples of $n$ constitute the principle ideal $[n]$, so the result means that $[m]$ is contained in $[n]$. Conversely $[m] \subseteq [n]$ means in particular that $m$ is in $[n]$, and hence $m= an$. Therefore 
\[ [m] \subseteq [n]  \mbox{ if and only if } n|m.\]
The ``bigger'' number corresponds to the ``smaller'' ideal; for instance the ideal $[6]$ of all multiples of 6 is properly contained in the ideal $[2]$ of all even integers. 


